\chapter{Producing, Consuming, Offsets, and Redelivery}
This chapter distinguishes four kinds of progress that are easy to confuse. A partition's log end offset (LEO) is the next offset available for append. During replication, the high watermark (HW) is a next-offset boundary; only records with offset < HW are within the replication-visible range. In the current single-broker path, HW equals LEO, and replication later in the book can separate them. A consumer's position is the in-memory offset for its next fetch; its committed next offset is progress stored per group and partition so it can be restored after restart. poll advances only the in-memory position, while commit writes progress to the offset log; neither deletes messages. A producer receipt reports the partition and half-open offset range \pathcode{[firstOffset,nextOffset)} for an append; it is not a consumer commit. A directly supplied \method{RpcClient} transfers to the client and closes with \method{close}; a shared \method{MetadataRouter} remains owned by its creator.

\lessonsection{13}{Batched Producer}
\goals{Implement per-partition ordered batch buffering, synchronous sends at the threshold, and return all receipts on explicit flush.}{Understand what a producer, partition routing, batch, acknowledgement, and unknown request outcome mean.}
\background{A producer is a client that writes records to a broker; each record is routed to a topic (a named category) and partition (one ordered log). A batch is a group of records temporarily held by the client for one send. Each partition buffers independently, so reaching the threshold for one partition does not send another partition's records. A receipt is the actual offset range returned by the broker; ACK is the requested acknowledgement mode. This step's single-broker backend ignores ACK; it also does not wait for a nonnegative \pathcode{ProduceRequest.timeoutMillis}. The producer waits synchronously for the broker reply, while the actual TCP wait limit is configured separately by \method{RpcClient}; this is not multi-replica confirmation. Replica ACK semantics arrive in \stepref{23}{produce acknowledgments}. An error or timeout does not prove that no append occurred: the request may have reached the broker and been written while its reply was lost. The course's ordinary \method{SimpleProducer} does not retry or use a background timer.}
\subsubsection*{Read first}
Review \stepref{3}{partitions and offsets within a partition}, then \stepref{9}{choosing a partition by key}, \stepref{10}{listing partitions through metadata}, and \stepref{12}{broker PRODUCE handling}. In particular, \method{Partitioner.choose} is student work completed in Step 9, not a provided answer; this step currently uses the direct \method{RpcClient} path. Client routing through \method{MetadataRouter} comes in Step 27.
\providededit{\method{RpcClient}, message/record value types, and the broker metadata/produce protocol are provided; reuse the Step 9 \method{Partitioner} implementation.}{\pathcode{exercises/src/main/java/io/simplekafka/client/SimpleProducer.java}: \method{void send(String topic, byte[] key, byte[] value, long timestamp)} and \method{List<ProduceReceipt> flush()}. Implement only these exercise methods; do not change \method{Partitioner}, the protocol, or tests.}
\subsubsection*{Inputs, outputs, and state}
\begin{itemize}
\item Constructor argument \pathcode{batchRecords} must be positive or reports \method{INVALID_REQUEST}; a null client/partitioner throws \pathcode{NullPointerException}. Default ACK is \pathcode{LEADER}, and the request field \pathcode{timeoutMillis} is 5,000 ms. Existing \method{setAcks(Acks,long)} changes ACK/timeout fields on later requests; timeout must be nonnegative, and policy changes are allowed only with no pending batch. Null acks throws \pathcode{NullPointerException}.
\item In \method{send}, topic must be a valid course identifier; key may be null; value must be non-null; timestamp is a nonnegative long, and the course does not require it to represent a real-world clock. A successful call returns no receipt: it appends the record to that \pathcode{TopicPartition}'s FIFO batch and synchronously sends that partition when the threshold is reached.
\item \method{flush()} sends every nonempty partition batch, returns an unmodifiable list of receipts accumulated since the previous flush, and clears that accumulation. With no pending records it sends no empty request and returns an empty list.
\item A metadata/parameter error before sending causes no broker append. An error after entering the produce RPC/reply path may follow an append; the producer saves the first error and enters an unknown-outcome state. Later \method{send}, \method{flush}, and \method{setAcks} calls report \method{REQUEST_TIMEOUT}; close the producer and create a new instance.
\end{itemize}
\lessoncontract{Preserve send order within a partition; different partitions have independent batches. Keep receipts from automatic sends until the next explicit flush. Receipts use half-open offset ranges. Null keys use the partitioner's per-instance rotation, and non-null keys use the Step 9 routing implementation.}
\subsubsection*{Separately checkable boundaries}
\begin{itemize}
\item A null value, negative timestamp, or invalid topic reports \method{INVALID_REQUEST} before entering produce; an unknown topic reports \method{UNKNOWN_TOPIC_OR_PARTITION}. These checks do not append or put the producer into an unknown-outcome state.
\item An oversized ordinary record (more than 1,048,580 bytes total; with a null key, the value limit is 1,048,548 bytes) reports \method{INVALID_REQUEST}; the broker receives no produce frame and appends nothing. If the batch threshold has not been reached, \method{send} may return successfully first; the local encoder discovers the size only on automatic send or \method{flush()}. At that point the producer marks the produce-path exception as an unknown outcome, and later operations report \method{REQUEST_TIMEOUT}. This is the course client's conservative state behavior, not evidence that the broker may have appended; do not resend the batch.
\item If any exception occurs on the first synchronous batch-send path, the call propagates its original error code and latches the producer's unknown-outcome state; later \method{send}, \method{flush}, and \method{setAcks} calls report \method{REQUEST_TIMEOUT}, and the batch must not be resent. Whether an append actually occurred depends on where the exception arose: a network/reply loss may follow an append, while a local encoding error proves no frame was sent but still latches this state in the course implementation.
\item Calling \method{setAcks} with buffered records or a negative timeout reports \method{INVALID_REQUEST}. Calling a producer operation after close reports \pathcode{IllegalStateException}.
\end{itemize}
\expected{Work it by hand: \pathcode{orders-0} starts at LEO=0, and the producer is configured with \pathcode{batchRecords=2}. Two records with the same key buffer in order; the second triggers the send, so the broker appends offsets 0 and 1 and returns \pathcode{[0,2)}. A third record remains buffered and is not yet readable. The next \method{flush()} sends it and returns two receipts, \pathcode{[0,2)} and \pathcode{[2,3)}; a second flush returns an empty list. Assume the topic has the target partition or the key is known to route there.\\
\method{Step13Test.flushesAtTheBatchThresholdAndReturnsAutomaticAndExplicitRanges} checks threshold behavior and one-time receipt delivery; \method{automaticFlushForOnePartitionLeavesAnotherPartitionsPendingBatchInvisible} checks independent partition batches; \method{rejectsInvalidSendsWithoutPersistingAnyRecords} and \method{doesNotRetryAProduceWhoseAppendSucceededBeforeItsResponseTimedOut} check pre-append rejection and no retry after append/reply loss.}
\lessonhints{First separate the times when a request is sent, the broker appends, the client receives a reply, and the caller receives a receipt; they need not coincide.}{Use \pathcode{TopicPartition} as the batch boundary; work through one partition reaching its threshold while another remains below it.}{Use the test case where the broker appends and then delays/drops its reply. Check whether the client preserves the first error and avoids sending the old batch again.}
\lessonquestions{Why can a timeout not prove that the record was not written?}{If an automatic send succeeds, why must the receipt still be returned by a later flush?}
\paragraph{Self-check explanation.}The reply can be lost after append, so the client cannot know the result; the receipt lets the caller identify the exact offset range, and dropping it makes an append impossible to track.
\pitfalls{Treating “received an error” as “definitely not written” and retrying may create duplicates. Real Kafka also has asynchronous sends and more acknowledgement choices; this course's \method{SimpleProducer} is synchronous, explicitly flushed, and does not retry.}
\lessoncommand{13}

\lessonsection{14}{Pull-Based Consumer}
\goals{Build a single-threaded, single-round pull consumer with explicit assignment and in-memory positions.}{Understand “next offset,” replay with seek, empty fetches, and the shared cross-partition budget.}
\background{A consumer actively pulls records from a broker; reading does not delete them from the log. Its position is the offset for the next fetch in one partition, not persisted progress. Manual assign establishes the set of partitions this client will read; seek changes only the local read position. poll issues one bounded fetch round rather than waiting continuously.}
\subsubsection*{Read first}
Review \stepref{3}{partitions and their local offsets} and \stepref{12}{broker FETCH requests}. Step 13 is the preceding course step, but this lesson does not call the producer or read OffsetStore. It starts with purely in-memory positions at offset 0; restoration and commit arrive in Step 15.
\providededit{The broker's \pathcode{FETCH} dispatch, \method{RpcClient}, and the \method{LogRecord} and \method{TopicPartition} value types are provided.}{\pathcode{exercises/src/main/java/io/simplekafka/client/SimpleConsumer.java}: \method{void assign(Collection<TopicPartition> partitions)}, \method{void seek(TopicPartition tp, long offset)}, and \method{Map<TopicPartition,List<LogRecord>> poll(int maxRecords, int maxBytes)}. Implement these three methods; the read-only \method{position(TopicPartition)} accessor is already provided.}
\subsubsection*{Inputs, outputs, and state}
\begin{itemize}
\item \method{assign} replaces the old assignment: deduplicate and sort by topic/partition; preserve positions for retained partitions, remove revoked ones, and initialize new ones at 0. A null collection throws \pathcode{NullPointerException}; a null element reports \method{INVALID_REQUEST}; rejection leaves the old assignment unchanged.
\item \method{seek(tp,offset)} requires an assigned partition and a nonnegative next offset. A null tp throws \pathcode{NullPointerException}; an unassigned partition reports \method{NOT_ASSIGNED}; a negative offset reports \method{OFFSET_OUT_OF_RANGE}; a rejected seek leaves positions unchanged. seek does not commit.
\item Both \pathcode{maxRecords} and \pathcode{maxBytes} in \method{poll} must be positive; otherwise report \method{INVALID_REQUEST} without advancing positions. They are total budgets shared across partitions, scanned in topic/partition order. An empty assignment sends no fetch; an empty partition spends no budget and is omitted from the returned map. A successful returned map is unmodifiable. After close, consumer operations throw \pathcode{IllegalStateException}; \method{position(tp)} on an unassigned partition reports \method{NOT_ASSIGNED}.
\item At this stage an ordinary record costs \pathcode{32 + keyBytes + valueBytes}; a record cannot be split. Each successful partition reply must be contiguous from its requested position and fit the remaining record/byte budgets. Otherwise that partition reports \method{INVALID_REQUEST} and does not advance. If seek starts before a deleted log start or after the current LEO, the broker reports \method{OFFSET_OUT_OF_RANGE} and the failed partition does not advance. Other server errors propagate unchanged; positions of earlier successful partitions may already have advanced while the failed partition remains at its old position.
\end{itemize}
\lessoncontract{A partition that returns records advances to the last returned offset plus one. A partition with no returned records, an empty fetch, or a failed fetch does not advance. poll does not commit or delete log data. The broker caps oversized positive maxBytes to a conservative single-frame budget; a later round continues at the first unreturned offset.}
\expected{Assume \pathcode{orders-0} and \pathcode{orders-1} each contain two ordinary records, each with null key and empty value, so each record costs 32 bytes; both positions start at 0. After assigning both partitions, poll(4,64) scans first \pathcode{orders-0}, returns offsets 0 and 1, and produces a map containing only that partition; positions become (2,0) and the budget is exhausted. Seek both back to 0, then poll(4,63): the first partition can return only offset 0, since the remaining 31 bytes cannot fit another record, so positions become (1,0). In a scripted test, if the broker replies \method{UNKNOWN_TOPIC_OR_PARTITION} for the second partition (simulating that request's failure), the first position may be 1 while the failed one stays 0.\\
\method{Step14Test.pollsInPartitionOrderWithinRoundBudgetsAndPreservesManualPositions} checks ordering, budgets, seek, and assign; \method{pollContinuesPastEmptyPartitionsWithoutSpendingTheSharedBudget} checks empty partitions; \method{leavesFailedPartitionPositionUnchangedForFetchErrorsAndMalformedSuccessReplies} checks failures, discontinuous offsets, and oversized replies.}
\lessonhints{Track assignment, each in-memory position, and persisted committed offset separately; this step changes only the first two.}{Subtract each whole record's encoded bytes and one record from the same poll's remaining budget; do not reset it per partition.}{Let an earlier-sorted partition succeed and a later one fail. Inspect both positions separately instead of assuming the whole round rolls back atomically.}
\lessonquestions{After poll reads offsets 0 and 1 and sets position to 2, why might a rebuilt consumer still read from 0?}{Why must a later partition's position not advance when its fetch fails?}
\paragraph{Self-check explanation.}The position lives only in the old client's memory; without resume/commit, a new assignment defaults to 0. A failed partition was not successfully read, so advancing it could skip data, while an earlier successful partition may retain its own progress.
\pitfalls{A successful read is not a durable commit; an empty read at the log end is no reason to increment position. Real Kafka has continuous polling, heartbeats, and a more complex assignment lifecycle; this step implements only single-threaded bounded rounds.}
\lessoncommand{14}

\lessonsection{15}{Consumer Offset Persistence}
\goals{Persist each group's next partition offset in a local durable log and restore it in a new consumer.}{Distinguish consumer progress from OffsetStore's physical log offsets and its group key.}
\background{A consumer position is an in-memory read cursor; commit stores the “next offset to read” as consumer progress. OffsetStore is itself an append log, but its physical offsets number the store records; the \pathcode{nextOffset} in each payload is the consumer progress. The key is \pathcode{(group,TopicPartition)}, so a different group or partition has separate progress. Reading or committing neither deletes business records nor changes HW.}
\subsubsection*{Read first}
You need \stepref{1}{log records and recovery}, \stepref{12}{broker request dispatch}, and \stepref{14}{position, assign, and poll}. A manual consumer uses a null group token; membership fencing does not arrive until Step 20.
\providededit{Offset-entry encoding in \method{MessageCodec}, the backing \method{PartitionLog}, record recovery, and \pathcode{COMMIT_OFFSET/FETCH_OFFSET} request types are provided.}{\pathcode{exercises/src/main/java/io/simplekafka/group/OffsetStore.java}: \method{void commit(OffsetKey key, long nextOffset)}, \method{OptionalLong fetch(OffsetKey key)}; \pathcode{exercises/src/main/java/io/simplekafka/client/SimpleConsumer.java}: \method{void resume(Collection<TopicPartition> partitions)}, \method{void commitSync()}; also extend \method{handle(Messages.Request)} in \pathcode{exercises/src/main/java/io/simplekafka/broker/BrokerHandler.java} with \pathcode{COMMIT_OFFSET/FETCH_OFFSET} branches.}
\subsubsection*{Inputs, outputs, and state}
\begin{itemize}
\item \pathcode{OffsetKey} identifies a group and one topic-partition. \method{commit}'s nextOffset must be nonnegative; moving it backward is allowed. A null commit key, negative value, or null fetch key reports \method{INVALID_REQUEST}; a closed store or storage failure reports \method{STORAGE_ERROR}. \method{fetch} returns the latest \pathcode{OptionalLong}; no commit is empty, not a committed offset of 0.
\item Each commit appends and forces an encoded offset entry with a null record key and timestamp 0. On reopen, the last record for each key wins. The store's physical log offset need not equal the payload's nextOffset. A complete corrupt record reports \method{CORRUPT_RECORD} without changing the file; a valid but incomplete tail may be truncated by the backing recovery code.
\item \method{resume} deduplicates and sorts partitions, replaces the assignment, and sets each position to its committed value or 0 if missing. It does not clamp a committed value to the current LEO. A lookup failure propagates after the assignment has changed and is not rolled back: new partitions remain at 0 and retained partitions keep their old positions. \method{commitSync} sends one request per assigned partition, including position 0; an empty assignment sends no commit request. poll does not commit, and commits for several partitions are not one atomic transaction.
\item Protocol success bodies: \pathcode{CommitOffsetRequest(OffsetKey,nextOffset,null)} returns \pathcode{EmptyBody}; \pathcode{FetchOffsetRequest(OffsetKey)} returns an \pathcode{OffsetBody} with an OptionalLong. A negative offset is also \method{INVALID_REQUEST} through the broker.
\end{itemize}
\lessoncontract{The most recently appended value for the same key is current, and it may move backward; changing either group or topic-partition isolates the value. resume restores durable progress but does not repair an out-of-range commit or write the consumer's current position to the store automatically.}
\expected{Start with no \pathcode{workers/orders-0} value. Commit 3, close and reopen, and fetch returns 3; commit 1 and reopen again, and fetch returns 1, demonstrating latest-wins and rewind. The same partition under another group is still empty. A new consumer resumes at position 1. If the stored value is 99 while the partition LEO is 3, resume still sets 99; poll later reports \method{OFFSET_OUT_OF_RANGE} and does not change the stored value.\\
\method{Step15Test.offsetStorePersistsNextOffsetsPerGroupAndPartitionAndAllowsRewind} checks isolation, rewind, and reopen; \method{rawOffsetRequestsReturnExactBodiesAndIsolateGroupTopicAndPartitionKeys} checks protocol bodies; \method{resumeCommitsAllPositionsAndDoesNotClampPastLogEnd} checks resume; \method{offsetStoreTruncatesAnIncompleteTailAndRejectsCompleteBadCrcWithoutMutation} checks recovery and complete corruption.}
\lessonhints{First treat consumer progress as the next business-record offset, not as the physical offset in OffsetStore's log.}{One OffsetKey can have several historical commits; use the final append, not an assumed map order.}{Test two separate facts: whether the value remains after reopening the store, and whether a fresh consumer.resume starts at that value.}
\lessonquestions{After processing offset 2, why commit nextOffset=3 instead of 2?}{If nextOffset is 99 and LEO is 3, should resume silently change it to 3?}
\paragraph{Self-check explanation.}nextOffset marks the next record to process, so 3 skips already processed offsets 0--2. The store restores the caller's committed value faithfully; a later fetch reports the out-of-range condition instead of silently rewriting progress.
\pitfalls{Committing the last processed offset makes it replay after recovery; confusing committed offset with a replication visibility boundary is also wrong. Real Kafka stores offsets in a replicated/compacted internal topic; this course uses a local single-coordinator append log and provides no high availability.}
\lessoncommand{15}

\lessonsection{16}{At-Least-Once Processing}
\goals{Order polling, business side effects, and offset commits safely, and work out replay after failure.}{Understand at-least-once delivery, in-memory positions versus durable commits, and non-atomic requests.}
\background{A business side effect changes state while processing a record, for example writing a local ledger or sending a notification. One call to \method{ProcessingLoop} polls first, invokes callbacks in ascending TopicPartition order and ascending offset within each partition, and calls \method{commitSync} only after every callback succeeds. Side effects and the offset log are not one transaction: if processing succeeds but commit fails, recovery must read the record again and may repeat the side effect. At-least-once means preferring a possible repeat of successful work over silently skipping work that did not succeed.}
\subsubsection*{Read first}
You need \stepref{14}{poll and position} and \stepref{15}{durable next offsets and commitSync}. The callback interface is provided; no transaction joining external effects to the offset store is required.
\providededit{The \method{RecordProcessor.process(TopicPartition,LogRecord)} callback interface and the one-round poll/commit lifecycle are provided.}{\pathcode{exercises/src/main/java/io/simplekafka/client/ProcessingLoop.java}: \method{int runOnce(int maxRecords, int maxBytes) throws Exception}. Implement the one-round controller only; do not change the consumer, store, or callback interface.}
\subsubsection*{Inputs, outputs, and state}
\begin{itemize}
\item maxRecords and maxBytes are positive total budgets across the poll's partitions; a nonpositive limit is rejected by poll with \method{INVALID_REQUEST}. On success return the count of records processed and committed. An empty poll still calls commitSync and returns 0.
\item The callback receives a partition and one \method{LogRecord}, returns no value, and may throw any \pathcode{Exception}. Its original exception propagates, and the round commits no partitions; poll may already have advanced the consumer's in-memory positions.
\item An exception from poll, callback, or commit marks this ProcessingLoop failed; a later runOnce throws \pathcode{IllegalStateException}. Close the old consumer, create a new consumer and loop, and resume from durable committed positions. If a side effect happened before its offset was committed, that effect may repeat.
\item After all callbacks succeed, \method{commitSync} commits each assigned partition separately; this is not a cross-partition transaction, so a later commit failure may follow successful earlier commits. The failure propagates and the method does not return a successful count.
\end{itemize}
\lessoncontract{Invoke callbacks in stable partition order and offset order; begin committing only after every record returned by this poll has been processed successfully. If any callback fails, do not commit any partition. After any failed round, do not reuse the loop or treat its advanced in-memory position as durable progress.}
\expected{Initially \pathcode{orders-0} and \pathcode{orders-1} each contain offsets 0 and 1, and group \pathcode{workers} has no commits. poll(4,8192) returns all four. Callbacks first cause effects \pathcode{p0:0,p0:1}, then throw at \pathcode{p1:0}; neither partition has a committed offset, although the old consumer's in-memory positions reached 2. Rebuild and resume: all four records are read from 0/0, so the two earlier p0 effects happen again. In a different failure, all callbacks finish but commit requests fail: effects already happened, while commits may be unconfirmed; recovery may replay, and earlier partition keys may already have committed.\\
\method{Step16Test.laterPartitionFailureLeavesBothOffsetsUncommittedAndReplaysEarlierSideEffects} covers a later-partition callback failure; \method{commitFailureAfterSideEffectsReplaysTheSameRecordsAtLeastOnce} covers commit failure; \method{boundedRoundsProcessInTopicPartitionOffsetOrderAndCommitFinalPositions} covers ordering, budgets, and empty rounds; \method{failedPollProcessesNothingAndResumeRetriesThePreviouslyFetchedRecord} covers poll failure.}
\lessonhints{Write poll, each side effect, each offset request, and crash/rebuild in temporal order.}{List callbacks in the required order, then decide whether failure happened in a callback or during commit; they have different side-effect histories.}{For each failure in Step16Test, record the store value, consumer position, and whether the loop can be reused separately; do not collapse them into one “progress” value.}
\lessonquestions{If p0 processing succeeds but p1 processing fails, why can p0's successful side effects also replay?}{Is commitSync all-or-nothing across three partitions?}
\paragraph{Self-check explanation.}The contract waits for all callbacks before beginning commits; if a callback fails, no partition is committed, so recovery reads from the old durable offsets. commitSync sends separate requests, not a cross-partition transaction, so an earlier partition may have committed already.
\pitfalls{Swallowing a callback exception and then committing can lose a failed record; continuing to use a failed loop can skip uncommitted data. At-least-once is not exactly-once. An application must make its side effects deduplicable; Step 31 implements local durable business deduplication.}
\lessoncommand{16}
